%%%%%%%%%%%%%%%%%%%%%%%%%%%%%%%%%%%%%%%%%%%%%%%%%%%%%%%%%%%%%%%%%%%%%%
% Lớp 8 | Chương II. Hằng đẳng thức đáng nhớ và ứng dụng | Hiệu hai bình phương. Bình phương của một tổng hay một hiệu
% Dạng 8D2?6-4: Chứng minh đẳng thức, tìm giá trị nhỏ nhất, lớn nhất
%%%%%%%%%%%%%%%%%%%%%%%%%%%%%%%%%%%%%%%%%%%%%%%%%%%%%%%%%%%%%%%%%%%%%%


% nguồn: dan_tu_AI (nhập 14/11/2026)
\begin{bt}%[8D2H6-4]%[ma:b759e795]
	Chứng minh các đẳng thức sau:\\
	a) $(a+b)^2-(a-b)^2=4ab$;\\
	b) $(a+b)^2+(a-b)^2=2\left(a^2+b^2\right)$.
	\loigiai{
		a) Biến đổi vế trái:
		\[ (a+b)^2-(a-b)^2=\left(a^2+2ab+b^2\right)-\left(a^2-2ab+b^2\right)=4ab. \]
		Vế trái bằng vế phải, đẳng thức được chứng minh.\\
		b) Biến đổi vế trái:
		\[ (a+b)^2+(a-b)^2=a^2+2ab+b^2+a^2-2ab+b^2=2a^2+2b^2=2\left(a^2+b^2\right). \]
		Vế trái bằng vế phải, đẳng thức được chứng minh.
	}
\end{bt}

% nguồn: dan_tu_AI (nhập 14/11/2026)
\begin{bt}%[8D2H6-4]%[ma:8f4ca59e]
	Chứng minh giá trị của biểu thức sau không phụ thuộc vào giá trị của biến $x$:
	\[ A=(2x-1)^2-(2x+3)(2x-3)+4x. \]
	\loigiai{
		Ta có
		\[ A=\left(4x^2-4x+1\right)-\left(4x^2-9\right)+4x=4x^2-4x+1-4x^2+9+4x=10. \]
		Vậy $A=10$ với mọi $x$, tức là giá trị của $A$ không phụ thuộc vào $x$.
	}
\end{bt}

% nguồn: dan_tu_AI (nhập 14/11/2026)
\begin{bt}%[8D2H6-4]%[ma:b003f431]
	Tìm giá trị nhỏ nhất của các biểu thức sau:\\
	a) $A=x^2-4x+7$;\qquad b) $B=x^2+x+1$.
	\loigiai{
		a) $A=x^2-2\cdot x\cdot 2+4+3=(x-2)^2+3$.\\
		Vì $(x-2)^2\ge 0$ với mọi $x$ nên $A\ge 3$; dấu "$=$" xảy ra khi $x-2=0$, tức $x=2$.\\
		Vậy giá trị nhỏ nhất của $A$ là $3$, đạt được khi $x=2$.\\
		b) $B=x^2+2\cdot x\cdot\dfrac12+\dfrac14+\dfrac34=\left(x+\dfrac12\right)^2+\dfrac34$.\\
		Vì $\left(x+\dfrac12\right)^2\ge 0$ với mọi $x$ nên $B\ge\dfrac34$; dấu "$=$" xảy ra khi $x=-\dfrac12$.\\
		Vậy giá trị nhỏ nhất của $B$ là $\dfrac34$, đạt được khi $x=-\dfrac12$.
	}
\end{bt}

% nguồn: dan_tu_AI (nhập 14/11/2026)
\begin{bt}%[8D2V6-4]%[ma:18414d58]
	a) Tìm giá trị lớn nhất của biểu thức $C=-x^2+6x-5$.\\
	b) Chứng minh $x^2-2x+2>0$ với mọi số thực $x$.
	\loigiai{
		a) $C=-\left(x^2-6x+9\right)+9-5=4-(x-3)^2$.\\
		Vì $(x-3)^2\ge 0$ với mọi $x$ nên $C\le 4$; dấu "$=$" xảy ra khi $x=3$.\\
		Vậy giá trị lớn nhất của $C$ là $4$, đạt được khi $x=3$.\\
		b) $x^2-2x+2=\left(x^2-2x+1\right)+1=(x-1)^2+1$.\\
		Vì $(x-1)^2\ge 0$ với mọi $x$ nên $(x-1)^2+1\ge 1>0$ với mọi $x$.
	}
\end{bt}

% nguồn: dan_tu_AI (nhập 14/11/2026)
\begin{bt}%[8D2V6-4]%[ma:3698e5d4]
	a) Tìm giá trị nhỏ nhất của biểu thức $P=x^2+y^2-4x+2y+10$.\\
	b) Bác Hùng có $20\text{ m}$ lưới để rào một mảnh vườn hình chữ nhật (rào kín bốn cạnh). Gọi $x$ (m) là chiều dài một cạnh của mảnh vườn ($0<x<10$). Viết biểu thức tính diện tích $S$ của mảnh vườn theo $x$ và tìm giá trị lớn nhất của $S$.
	\loigiai{
		a) $P=\left(x^2-4x+4\right)+\left(y^2+2y+1\right)+5=(x-2)^2+(y+1)^2+5$.\\
		Vì $(x-2)^2\ge 0$, $(y+1)^2\ge 0$ nên $P\ge 5$; dấu "$=$" xảy ra khi $x=2$ và $y=-1$.\\
		Vậy giá trị nhỏ nhất của $P$ là $5$, đạt được khi $x=2$, $y=-1$.\\
		b) Nửa chu vi mảnh vườn là $10\text{ m}$ nên cạnh còn lại dài $10-x$ (m).\\
		\[ S=x(10-x)=10x-x^2=-\left(x^2-10x+25\right)+25=25-(x-5)^2. \]
		Vì $(x-5)^2\ge 0$ nên $S\le 25$; dấu "$=$" xảy ra khi $x=5$.\\
		Vậy diện tích lớn nhất là $25\text{ m}^2$, khi mảnh vườn là hình vuông cạnh $5\text{ m}$.
	}
\end{bt}

% nguồn: De_AI_Toan8_cuoi_hoc_ki_i_1 (nhập 21/11/2026)
\begin{bt}%[8D2N6-4]%[ma:c7b7a3d7]
	Chứng minh rằng giá trị của biểu thức $K = (x + 5)^2 - x(x + 10)$ không phụ thuộc vào giá trị của biến $x$.
	\loigiai{
		Ta có $K = (x + 5)^2 - x(x + 10) = (x^2 + 10x + 25) - (x^2 + 10x) = x^2 + 10x + 25 - x^2 - 10x = 25$. \\
Vì $K = 25$ là một hằng số nên giá trị của $K$ không phụ thuộc vào giá trị của biến $x$.
	}
\end{bt}

% nguồn: De_AI_Toan8_cuoi_hoc_ki_i_1 (nhập 21/11/2026)
\begin{bt}%[8D2N6-4]%[ma:59445e4a]
	Chứng minh các đẳng thức sau:
a) $(a + b)^2 - (a - b)^2 = 4ab$;
b) $(a + b)^2 + (a - b)^2 = 2(a^2 + b^2)$.
	\loigiai{
		a) Biến đổi vế trái: $(a + b)^2 - (a - b)^2 = (a^2 + 2ab + b^2) - (a^2 - 2ab + b^2) = a^2 + 2ab + b^2 - a^2 + 2ab - b^2 = 4ab$ (bằng vế phải). Điều phải chứng minh. \\
b) Biến đổi vế trái: $(a + b)^2 + (a - b)^2 = (a^2 + 2ab + b^2) + (a^2 - 2ab + b^2) = 2a^2 + 2b^2 = 2(a^2 + b^2)$ (bằng vế phải). Điều phải chứng minh.
	}
\end{bt}

% nguồn: De_AI_Toan8_cuoi_hoc_ki_i_1 (nhập 21/11/2026)
\begin{bt}%[8D2N6-4]%[ma:77dbe699]
	Chứng minh rằng biểu thức $P = x^2 - 6x + 14$ luôn nhận giá trị dương với mọi giá trị của $x$, từ đó tìm giá trị nhỏ nhất của $P$.
	\loigiai{
		Ta có $P = x^2 - 6x + 14 = (x^2 - 6x + 9) + 5 = (x - 3)^2 + 5$. \\
Vì $(x - 3)^2 \ge 0$ với mọi $x$, nên $P = (x - 3)^2 + 5 \ge 5 > 0$ với mọi $x$. \\
Dấu bằng xảy ra khi $x - 3 = 0 \Leftrightarrow x = 3$. \\
Vậy giá trị nhỏ nhất của $P$ là $5$, đạt được khi $x = 3$.
	}
\end{bt}
